\documentclass[10pt,a4paper]{amsart}
\usepackage[margin=19mm]{geometry}
\usepackage{amsmath,amssymb,mathtools}
\usepackage[scr=rsfs,cal=euler]{mathalfa}
\usepackage{xfrac}
\usepackage[unicode,hidelinks,bookmarks=false]{hyperref}
\newcommand{\ie}{i.e.\ }
\newcommand{\cf}{cf.\ }
\newcommand{\resp}{respectively\ }
\newcommand{\Subsection}{\subsection}
\providecommand{\nobreakdash}{\nobreak\hbox{-}}
\setlength{\emergencystretch}{2em}
\multlinegap0pt
\usepackage{amssymb}
\usepackage{mathtools}
\usepackage{xcolor}
\newtheorem{definition}{Definition}
\newtheorem{proposition}{Proposition}
\newtheorem{lemma}{Lemma}
\newtheorem{corollary}{Corollary}
\newtheorem{remark}{Remark}
\newcommand{\N}{\mathbb{N}}
\newcommand{\R}{\mathbb{R}}
\newcommand{\dotcirc}{\mathbin{\odot}}
\newcommand{\e}{\mathrm{e}}
\renewcommand{\epsilon}{\varepsilon}
\newcommand{\la}{\langle}
\newcommand{\noi}{\noindent}
\newcommand{\ra}{\rangle}
\title{Gauge transforms, random averaging operator ansatz and improved probabilistic well-posedness for the radial NLS on the $3d$ ball}
\author{Nicolas Burq, Nicolas Camps, Chenmin Sun, Nikolay Tzvetkov}
\date{Source: arXiv:2606.07010v1 (CC BY 4.0)}
\begin{document}
\maketitle

\noindent\emph{Source and licence.} Gauge transforms, random averaging operator ansatz and improved probabilistic well-posedness for the radial NLS on the $3d$ ball. Version arXiv:2606.07010v1 (CC BY 4.0), distributed by arXiv under the Creative Commons Attribution 4.0 International licence, http://creativecommons.org/licenses/by/4.0/, as recorded in the arXiv licence metadata for this exact version and shown on the abstract page https://arxiv.org/abs/2606.07010v1. Full licence notice: \texttt{licenses/arxiv-2606.07010-CC-BY-4.0.txt}.\par\smallskip

\noindent\textit{Editorial scope.} Counted unit: the article's proposition prop:quintilinear (source0.tex lines 1138--1164). The article's complete original proof of it is delivered with it (source0.tex lines 4745--5011, one \texttt{proof} environment of 267 lines, which the article places away from the statement). No mathematics is written, repaired or re-derived: the statement and the proof are the article's own lines, in the article's own notation, and no retained line is substituted or re-flowed.  Retained uncounted as the source-local context the proof needs: lines 754--763; lines 920--945; lines 3239--3252; lines 3517--3554; lines 3678--3701; lines 3703--3712; lines 3709--3711; lines 3853--3864; lines 4643--4663.  Nothing was omitted from any retained span, so the package carries no omission record.  No retained line contains a citation, so the wrapper carries no bibliography and no bibliography entry.  No retained line contains a figure environment, an image-inclusion command or drawing code, so no image is delivered and the package declares no asset dependency.  Editorial formatting only, in addition: an AMS \texttt{amsart} preamble that restates the article's own declarations and macros used by the retained lines, namely the environments \texttt{definition}, \texttt{proposition}, \texttt{lemma}, \texttt{corollary}, \texttt{remark} on separate counters, together with the standard packages those lines need.  The retained environments print in this document as Definition 1, Proposition 1, Lemma 1, Corollary 1, Remark 1 in order of appearance, whereas the article prints the same items with the numbering of its own full document.  The complete native source of the article as distributed in the arXiv e-print for this exact version is bundled unchanged as \texttt{sources/202225/source0.tex}.
\begin{definition}[Hierarchy of small parameters]\label{hierarchy}
	Fix a universal small parameter~$0<\sigma\ll 1$, and set~$(s,q,\gamma,b,\delta)=(s_{\sigma},q_{\sigma},\gamma_{\sigma},b_{\sigma},\delta_{\sigma})$ satisfying
	\begin{align*}
		&\frac{1}{q_{\sigma}}=\sigma, \qquad
		b_{\sigma}-\frac{1}{2}=\gamma_{\sigma}-\frac{1}{q_{\sigma}'}=\sigma^{10},\qquad
		\delta_{\sigma}=\sigma^{20},\\
		&\frac{1}{2}+2^{200}\sigma<s_{\sigma}<4\alpha-3-2^{100}\sigma,\qquad
		s_{\sigma}+4\alpha-\frac{9}{2}>2^{100}\sigma.
	\end{align*}
\end{definition}

\begin{definition}\label{def:CD}
	Assume that~$(s,q,\gamma,b,\delta)=(s_{\sigma},q_{\sigma},\gamma_{\sigma},b_{\sigma},\delta_{\sigma})$ satisfies the hierarchy in Definition~\ref{hierarchy}.
	We define type~$\mathrm{(C)}$ and type~$\mathrm{(D)}$ terms as follows:
	\begin{enumerate}
		\item Type~$\mathrm{(C)}$: an input function~$v$ is of type~$\mathrm{(C)}$ at scale~$N$ of size~$R$, if there is a~$\mathcal{B}_{\leq \frac{N}{2}}$-measurable function~$\mathcal{T}_n^N(t)\in\widetilde{\mathcal{F}}L_q^{\gamma}(\R)$, compactly supported in~$t\in(-2T,2T)$, such that
		\[
		v(t,x) = \sum_{\frac{N}{2}<n\leq N}
		\mathcal{T}^N_{n}(t)\frac{g_{n}(\omega)}{n^{\alpha}}\e^{-it\lambda_{n}^{2}}\cdot\mathbf{e}_{n}(x)\,,
		\]
		satisfying
		\begin{align}\label{bddef:TypeC}
		\sup_{\frac{N}{2}<n\leq N}\big(\|\langle\kappa\rangle^{\gamma}\widetilde{\mathcal{T}}^N_{n}(\kappa)\|_{L_{\kappa}^{q}} + \|\la\kappa\ra^b\widetilde{\mathcal{T}}^N_n(\kappa) \|_{L_{\kappa}^2}\big)
		\leq R  \,,
		\end{align}
		and
		\begin{align}\label{bddef:TypeC'}
		&\|v(t)\|_{X^{0,b}}\leq RT^{-(b-\frac{1}{2})}N^{-\alpha+\frac{1}{2}+\delta}, \notag \\ &\|v(t)\|_{X_{q,q}^{0,\gamma}}\leq RT^{-(\gamma-\frac{1}{q'})}N^{-\alpha+\frac{1}{q}+\delta}.
		\end{align}
		\item Type~$\mathrm{(D)}$: an input function~$v$ is of type~$\mathrm{(D)}$ at scale~$N$, if~$v$ is compactly supported in~$t\in(-T,T)$, with the estimates
		\[
		\|v\|_{X^{0,b}}
		\leq N^{-s},\quad \|(\mathrm{Id}-\Pi_L)v\|_{X^{0,b}}\leq  \big(\frac{N}{L}\big)N^{-s}
		\]
		for~$L\geq 4N$.
	\end{enumerate}
\end{definition}

\begin{lemma}\label{lem:N4mapping}
	Let $q_0\in[2,\infty)$, $\gamma_0\in(\frac{1}{q_0'},1)$, $0<\delta_0<\frac{1-\gamma_0}{2}$, and $r_0=\big(\gamma_0-\frac{1}{q_0'}+2\delta_0\big)^{-1}$. Then for any $F\in X_{\infty,r_0}^{0,0}$, $u\in X_{q_0,q_0}^{0,\gamma_0}$, and $0<T<1$,
	\[
	\|\chi_T(t)\,\mathcal{I}(F\dotcirc u)_n(t)\|_{\widetilde{\mathcal{F}}L_{q_0}^{\gamma_0}}
	\lesssim
	T^{\delta_0+\frac{1}{r_0'}}\|\chi_1(t)F_n(t)\|_{L_t^{\infty}}\|u_n\|_{\widetilde{\mathcal{F}}L_{q_0}^{\gamma_0}}.
	\]
	Consequently,
	\[
	\|\chi_T(t)\,\mathcal{I}(F\dotcirc u)\|_{X_{q_0,q_0}^{0,\gamma_0}}
	\lesssim
	T^{\delta_0+\frac{1}{r_0'}} \|u\|_{X_{q_0,q_0}^{0,\gamma_0}}\|F\|_{\ell_n^{\infty}L_t^{\infty}}.
	\]
\end{lemma}

\begin{corollary}\label{cor:4-linear-1}
	Assume $\alpha>15/16$ and $(s,q,\gamma,b,\delta)=(s_{\sigma},q_{\sigma},\gamma_{\sigma},b_{\sigma},\delta_{\sigma})$ satisfies the hierarchy in Definition~\ref{hierarchy}, with $\theta=b-\frac{1}{2}=\gamma-\frac{1}{q'}$. Assume moreover $M_{(1)}\sim M_{(2)}$, and let $M=M_{(1)}$, and assume that $\mathbf{P}_{M_j}v_{M_j}^{(\ast_j)}=v_{M_j}^{(\ast_j)}$ for $j=1,2,3,4$.
	Then the following bounds hold:
	\begin{align*}
		\mathrm{(1)} \quad &
		\|\chi(t)\,\mathcal{N}_n^{(4)}(v^{(\mathrm{D})}_{M_1},v^{(\mathrm{D})}_{M_2},v^{(\mathrm{D})}_{M_3},v^{(\mathrm{D})}_{M_4})\|_{L_t^{\infty}}
		\lesssim
		R^4 M_{(1)}^{10\theta+\delta}\cdot
		\frac{M_{(1)}\wedge n}{M_{(1)}}\cdot M_{(1)}^{-(3\alpha+2s-\frac{7}{2})},\\
		\mathrm{(2)} \quad &
		\|\chi(t)\,\mathcal{N}_n^{(4)}(v^{(\mathrm{D})}_{M_1},v^{(\mathrm{D})}_{M_2},v^{(\mathrm{D})}_{M_3},v^{(\mathrm{C})}_{M_4})\|_{L_t^{\infty}}
		\lesssim
		R^3 M_{(1)}^{10\theta+\delta}\cdot
		\frac{M_{(1)}\wedge n}{M_{(1)}}\cdot M_{(1)}^{-(3\alpha+2s-\frac72)},\\
		\mathrm{(3)} \quad &
		\|\chi(t)\,\mathcal{N}_n^{(4)}(v^{(\mathrm{D})}_{M_1},v^{(\mathrm{D})}_{M_2},v^{(\mathrm{C})}_{M_3},v^{(\mathrm{D})}_{M_4})\|_{L_t^{\infty}}
		\lesssim
		R^3 M_{(1)}^{10\theta+\delta}\cdot
		\frac{M_{(1)}\wedge n}{M_{(1)}}\cdot M_{(1)}^{-(3\alpha+2s-\frac72)},\\
		\mathrm{(4)} \quad &
		\|\chi(t)\,\mathcal{N}_n^{(4)}(v^{(\mathrm{D})}_{M_1},v^{(\mathrm{D})}_{M_2},v^{(\mathrm{C})}_{M_3},v^{(\mathrm{C})}_{M_4})\|_{L_t^{\infty}}
		\lesssim
		R^2 M_{(1)}^{10\theta+\delta}\cdot
		\frac{M_{(1)}\wedge n}{M_{(1)}}\cdot M_{(1)}^{-(3\alpha+2s-4)},\\
		\mathrm{(5)} \quad &
		\|\chi(t)\,\mathcal{N}_n^{(4)}(v^{(\mathrm{C})}_{M_1},v^{(\mathrm{C})}_{M_2},v^{(\mathrm{D})}_{M_3},v^{(\mathrm{D})}_{M_4})\|_{L_t^{\infty}}
		\lesssim
		R^2 M_{(1)}^{10\theta+\delta}\cdot
		\frac{M_{(1)}\wedge n}{M_{(1)}}\cdot M_{(1)}^{-(3\alpha+2s-4)}.
	\end{align*}
	In particular, for all interactions with at least two type~(D) inputs,
	\begin{align}\label{cor:4-linear-1-uniform}
		\|\chi(t)\,\mathcal{N}_n^{(4)}(v^{(\ast_1)}_{M_1},v^{(\ast_2)}_{M_2},v^{(\ast_3)}_{M_3},v^{(\ast_4)}_{M_4})\|_{L_t^{\infty}}
		\lesssim
		R^4 M_{(1)}^{10\theta+\delta}\cdot
		\frac{M_{(1)}\wedge n}{M_{(1)}}\cdot M_{(1)}^{-(3\alpha+2s-4)}.
	\end{align}
\end{corollary}

\begin{corollary}\label{cor:4-linearprob}
	Assume $\alpha>15/16$ and $(s,q,\gamma,b,\delta)=(s_{\sigma},q_{\sigma},\gamma_{\sigma},b_{\sigma},\delta_{\sigma})$ satisfies the hierarchy in Definition~\ref{hierarchy}, with $\theta=b-\frac{1}{2}=\gamma-\frac{1}{q'}$. Assume moreover $M_{(1)}\sim M_{(2)}$, and let $M=M_{(1)}$.
Then $(M^{\delta}R,\delta;\mathcal{B}_{\leq M})$-certainly,
	\begin{align*}
		\mathrm{(1)} \quad &
		\|\chi(t)\,\mathcal{N}_n^{(4)}(v^{(\mathrm{C})}_{M_1},v^{(\mathrm{C})}_{M_2},v^{(\mathrm{C})}_{M_3},v^{(\mathrm{D})}_{M_4})\|_{L_t^{\infty}}
		\lesssim
		R^3 M_{(1)}^{10\theta+\delta}\cdot
		\frac{M_{(1)}\wedge n}{M_{(1)}}\cdot M_{(1)}^{-(3\alpha+s-3)},\\
		\mathrm{(2)} \quad &
		\|\chi(t)\,\mathcal{N}_n^{(4)}(v^{(\mathrm{C})}_{M_1},v^{(\mathrm{C})}_{M_2},v^{(\mathrm{D})}_{M_3},v^{(\mathrm{C})}_{M_4})\|_{L_t^{\infty}}
		\lesssim
		R^3 M_{(1)}^{10\theta+\delta}\cdot
		\frac{M_{(1)}\wedge n}{M_{(1)}}\cdot
		M_{(1)}^{-(4\alpha+s-4)},\\
		\mathrm{(3)}\quad &
		\|\chi(t)\,\mathcal{N}_n^{(4)}(v^{(\mathrm{C})}_{M_1},v^{(\mathrm{C})}_{M_2},v^{(\mathrm{C})}_{M_3},v^{(\mathrm{C})}_{M_4})\|_{L_t^{\infty}}
		\lesssim
		R^4
		M_{(1)}^{10\theta+\delta}\cdot
		\frac{M_{(1)}\wedge n}{M_{(1)}}\cdot
		M_{(1)}^{-(5\alpha-4)}.
	\end{align*}
\end{corollary}

\begin{remark}\label{rmk:numerologyquintic}
	The hierarchy of parameters gives
	\[
	3\alpha+2s-4<s+4\alpha-4< \min\{s+3\alpha-3,\, 5\alpha-4\}.
	\]
	Combining this with Corollary~\ref{cor:4-linear-1}, we see that $(M^{\delta}R,\delta;\mathcal{B}_{\leq M})$-certainly, bounds~(1)--(5) in Corollary~\ref{cor:4-linear-1} together with~(1)--(3) of Corollary~\ref{cor:4-linearprob} are all controlled by the uniform estimate
	\begin{align}\label{finalbd:Nn4}
		R^4 M_{(1)}^{10\theta+\delta}\cdot\frac{M_{(1)}\wedge n}{M_{(1)}}\cdot M_{(1)}^{-(3\alpha+2s-4)}.
	\end{align}
\end{remark}

	\begin{align}
		R^4 M_{(1)}^{10\theta+\delta}\cdot\frac{M_{(1)}\wedge n}{M_{(1)}}\cdot M_{(1)}^{-(3\alpha+2s-4)}.
	\end{align}

\begin{definition}[Scale-$N$ $h^{s}$-envelope]\label{def:scale-seq}
	Let $\mathbf{a}=(a_M)_{M\in 2^{\mathbb N}}$ be a dyadic sequence, and for $\sigma\in\mathbb R$ define
	\[
	\|\mathbf{a}\|_{h^{\sigma}} := \bigl\| M^{\sigma} a_M \bigr\|_{\ell^2_{M\in 2^{\mathbb N}}}.
	\]
	(So $h^0=\ell^2$.) Fix $0<s\leq 1$. For a dyadic number $N\in 2^{\mathbb N}$, we say that $\mathbf{a}$ is a \emph{scale-$N$ $h^{s}$-envelope} if:
	\begin{align*}
		\text{(1)}\quad & \|\mathbf{a}\|_{h^0} \lesssim N^{-s},\\
		\text{(2)}\quad & \|\mathbf{a}\|_{h^{s_0}} \lesssim_{s_0,s} N^{-(s-s_0)},\qquad \forall\, s_0\in(0,s),\\
		\text{(3)}\quad & |a_M|\lesssim \Big(\frac{N}{M}\Big) N^{-s},\qquad \forall\, M\ge N.
	\end{align*}
\end{definition}

\begin{lemma}[Dyadic summation test]\label{lem:sum-abstract}
	Let $0<s_1,s_2,s_3,s_4\leq 1$. Assume that there exists
	\[
	\nu_0\in \Big(0,\min\{s_1,s_2,s_3,s_4\}\Big]
	\]
	such that for every dyadic quadruple $(M_1,M_2,M_3,M_4)\in (2^{\N})^4$ and every choice of scale-$M_j$ $h^{s_j}$-envelopes $b^{(j)}$, the typical-atom estimate
	\begin{equation}\label{eq:typical-atom-output-gain}
		c_{M_1,M_2,M_3,M_4}\prod_{j=1}^4 |b^{(j)}_{M_j}|
		\leq
		\Lambda\cdot M_{(1)}^{-\nu_0}
	\end{equation}
	holds for some $\Lambda>0$ independent of $M_1,M_2,M_3,M_4$.

	Then for any $(N_1,N_2,N_3,N_4)\in(2^{\N})^4$ and any scale-$N_j$ $h^{s_j}$-envelopes~$a^{(j)}$,
	\begin{equation}\label{eq:dyadic-sum-output-gain}
		\big|\mathcal{L}^{(4)}((a^{(j)})_{j=1}^4)\big|
		\lesssim
		\Lambda\cdot N_{(1)}^{-\nu_0}(\log N_{(1)})^4,
	\end{equation}
	where the implicit constant is independent of $N_j$.
\end{lemma}

\begin{proposition}[Quintilinear estimates]
	\label{prop:quintilinear}
 Assume that $R\geq 1$ and $T>0$. Let $(v_M^{(\mathrm{C})}, v_M^{(\mathrm{D})})_{M\in 2^{\N}}$ are sequences of type $(\mathrm{C})$ and type $(\mathrm{D})$ terms according to Definition \ref{def:CD}.
Then with an absolute constant $c_0>0$ and for all sufficiently small parameter $\sigma>0$ defining the parameters $(s,q,\gamma,b,\delta)$, the following estimates hold
$((2N)^{\delta}R,c_0;\mathcal{B}_{\leq2N})$-certainly for any choice of types~$(\ast_1,\ast_2,\ast_3)\in\{\mathrm{C},\mathrm{D}\}^3$:
	\begin{align*}
		\big\|\mathcal{I}_{\chi_T}\big(
		\mathcal{N}^{(4)}(v_{N_1}^{(\ast_1)},v_{N_2}^{(\ast_2)},v_{N_3}^{(\ast_3)},v_{N_4}^{(\ast_4)})
		\dotcirc v_{N_5}^{(\ast_5)}
		\big)\big\|_{X^{0,b}}
		&\lesssim R^5T^{\sigma}N_{(1)}^{-s}N_{(2)}^{-\delta},
	\end{align*}
	whenever~$N_5\leq \max(N_1,N_2,N_3,N_4)$, or~$2\max(N_1,N_2,N_3,N_4)\leq N_5$ and~$v_{N_5}^{(\ast_5)}$ is of type~(D).
	Moreover, if~$\max(N_1,N_2,N_3,N_4)\leq N$, the following operator bounds hold:
	\begin{align*}
		\big\|\cdot\mapsto
		\mathcal{I}_{\chi_T}\big(
		\cdot\dotcirc\mathcal{N}^{(4)}(v_{N_1}^{(\ast_1)},v_{N_2}^{(\ast_2)},v_{N_3}^{(\ast_3)},v_{N_4}^{(\ast_4)})
		\big)\big\|_{X^{0,b}\to X^{0,b}}
		&\lesssim R^4T^{\sigma}\max(N_1,\dots,N_4)^{-\delta},\\
		\big\|\cdot\mapsto
		\mathcal{I}_{\chi_T}\big(
		\cdot\dotcirc\mathcal{N}^{(4)}(v_{N_1}^{(\ast_1)},v_{N_2}^{(\ast_2)},v_{N_3}^{(\ast_3)},v_{N_4}^{(\ast_4)})
		\big)\big\|_{X_{q,q}^{0,\gamma}\to X_{q,q}^{0,\gamma}}
		&\lesssim R^4T^{\sigma}\max(N_1,\dots,N_4)^{-\delta}.
	\end{align*}
\end{proposition}

\begin{proof}[Proof of Proposition~\ref{prop:quintilinear}]

For $v_{N_j}^{(\ast_j)}$ with $\ast_j\in\{\mathrm{C},\mathrm{D}\}$ and $j\in\{1,2,3,4,5\}$, denote
\[
F_n(t):=\mathcal{N}_n^{(4)}(v_{N_1}^{(\ast_1)},v_{N_2}^{(\ast_2)},v_{N_3}^{(\ast_3)},v_{N_4}^{(\ast_4)})(t).
\]

By Lemma~\ref{lem:N4mapping},
\begin{align*}
	\big\|\chi_T(t)\,\mathcal{I}(F\dotcirc v_{N_5}^{(\ast_5)})_n\big\|_{\widetilde{\mathcal{F}}L_2^{b}}\lesssim T^{\delta_0+\frac{1}{r_0'}}\|\chi_1(t)F_n(t)\|_{L_t^{\infty}}\|v_{N_5,n}^{(\ast_5)}\|_{\widetilde{\mathcal{F}}L_2^b}.
\end{align*}
Assume that at least three of $v_{N_j}^{(\ast_j)}$, $j\in\{1,2,3,4\}$, are of type~(D). We apply Lemma~\ref{lem:sum-abstract} together with Corollary~\ref{cor:4-linear-1}.
In the proof of Corollary~\ref{cor:4-linear-1}, if $v^{(\ast_j)}$ is of type~(D), we use $\mathbf{a}^{(j)}=(\|\mathbf{P}_{M}v^{(\ast_j)}\|_{X^{0,b}})_{M\in 2^{\N}}$ as an $h^s$-envelope, while
for $v^{(\ast_j)}$ of type~(C), we use the $h^{s_C}$-envelope $\mathbf{b}^{(j)}=(\|\mathbf{P}_M v^{(\ast_j)}\|_{X_{q,q}^{0,\gamma}})_{M\in 2^{\N}}$, where
\[
s_C=\alpha-\frac{1}{q}-\delta,
\]
according to the different types of interactions in~(1)--(5) of Corollary~\ref{cor:4-linear-1}. In particular, the uniform estimate~\eqref{cor:4-linear-1-uniform} implies that for any given scales $M_1,M_2,M_3,M_4$, the typical-atom estimate holds with the bound $\Lambda M_{(1)}^{-\nu_0}$, where
\begin{align}\label{eq:Lambdanu_0}
	\Lambda:=R^4 T^{-3(\gamma-\frac{1}{q'})},\quad \nu_0=3\alpha+2s-4-10\theta-\frac{2}{q},
\end{align}
and one verifies from the hierarchy in Definition~\ref{hierarchy} that $\nu_0\leq \min\{s_C,s\}$. Applying Lemma~\ref{lem:sum-abstract}, we obtain
\begin{align}\label{bd:finalN_n^4}
	\|\chi_1(t)F_n(t)\|_{L_t^{\infty}}\lesssim_{\epsilon} \Lambda N^{-\nu_0+\epsilon},
\end{align}
where $N=\max(N_1,N_2,N_3,N_4)$.

By Lemma~\ref{lem:N4mapping}, we deduce
\begin{align}
	\big\|\cdot\mapsto
	\mathcal{I}_{\chi_T}\big(
	\cdot\dotcirc\mathcal{N}^{(4)}(v_{N_1}^{(\ast_1)},v_{N_2}^{(\ast_2)},v_{N_3}^{(\ast_3)},v_{N_4}^{(\ast_4)})
	\big)\big\|_{X^{0,b}\to X^{0,b}}
	&\lesssim R^4 T^{\delta+\frac{1}{r_0'}-3(\gamma-\frac{1}{q'})}N^{-\nu_0},\label{opquintic:X0b}\\
	\big\|\cdot\mapsto
	\mathcal{I}_{\chi_T}\big(
	\cdot\dotcirc\mathcal{N}^{(4)}(v_{N_1}^{(\ast_1)},v_{N_2}^{(\ast_2)},v_{N_3}^{(\ast_3)},v_{N_4}^{(\ast_4)})
	\big)\big\|_{X_{q,q}^{0,\gamma}\to X_{q,q}^{0,\gamma}}
	&\lesssim R^4 T^{\delta+\frac{1}{r_0'}-3(\gamma-\frac{1}{q'})}N^{-\nu_0}, \label{opquintic:X0q}
\end{align}
where $r_0^{-1}=\gamma-\frac{1}{q'}+2\delta=b-\frac{1}{2}+2\delta$. Under the hierarchy~\eqref{hierarchy},
\[
\delta+\frac{1}{r_0'}-3\Big(\gamma-\frac{1}{q'}\Big)\geq \sigma,
\quad\nu_0\geq \delta,
\]
which proves the last two inequalities of Proposition~\ref{prop:quintilinear}.

For the first inequality, we distinguish two cases.

\noi
$\bullet${\bf Case 1: $N_5=2N$ and $\max(N_1,N_2,N_3,N_4)\leq N$.}

We conclude by the operator norm estimate~\eqref{opquintic:X0b} and $\|v_{N_5}^{(\mathrm{D})}\|_{X^{0,b}}\leq (2N)^{-s}$.

\medskip

\noi
$\bullet${\bf Case 2: $N_5 \leq \max(N_1,N_2,N_3,N_4)=2N$.}

By Lemma~\ref{lem:N4mapping} and~\eqref{finalbd:Nn4},
\begin{align*}
	&\big\|
	\mathcal{I}_{\chi_T}\big(\mathcal{N}^{(4)}(v_{N_1}^{(\ast_1)},\cdots, v_{N_4}^{(\ast_4)})\dotcirc v_{N_5}^{(\ast_5)}\big)
	\big\|_{X^{0,b}}\\ &\qquad\leq
	\Big(\sum_{M}
	\big\|\mathcal{I}_{\chi_T}\big(\mathbf{P}_{M}\mathcal{N}^{(4)}(v_{N_1}^{(\ast_1)},\cdots, v_{N_4}^{(\ast_4)})\dotcirc \mathbf{P}_{M}v_{N_5}^{(\ast_5)}\big)
	\big\|_{X^{0,b}}^2
	\Big)^{1/2}\\
	&\qquad\lesssim T^{\delta+\frac{1}{r_0'}}\mathcal{Q},
\end{align*}
where
\begin{align}\label{eq:Q}
	\mathcal{Q}:=\Big(
	\sum_{M}\|\chi(t)\,\mathbf{P}_M\mathcal{N}_n^{(4)}(v_{N_1}^{(\ast_1)},\cdots, v_{N_4}^{(\ast_4)})\|_{\ell_n^{\infty}L_t^{\infty}}^2 \|\mathbf{P}_M v_{N_5}^{(\ast_5)}\|_{X^{0,b}}^2
	\Big)^{1/2}.
\end{align}


For $j\in\{1,2,3,4\}$, write
\[
v_{N_j}^{(\ast_j)}
=
\sum_{K_j\in 2^{\N}}\mathbf{P}_{K_j}v_{N_j}^{(\ast_j)}.
\]
For each fixed dyadic quadruple $(K_1,K_2,K_3,K_4)$, set $K:=\max(K_1,K_2,K_3,K_4)$.
For the output block, define
\[
b_M:=\|\mathbf{P}_M v_{N_5}^{(\ast_5)}\|_{X^{0,b}}.
\]
For the four input blocks, define
\[
a_{K_j}^{(j)}
:=
\begin{cases}
	\|\mathbf{P}_{K_j}v_{N_j}^{(\ast_j)}\|_{X^{0,b}},
	& v_{N_j}^{(\ast_j)} \text{ is of type (D)},\\[1mm]
	\|\mathbf{P}_{K_j}v_{N_j}^{(\ast_j)}\|_{X_{q,q}^{0,\gamma}},
	& v_{N_j}^{(\ast_j)} \text{ is of type (C)}.
\end{cases}
\]
By Definition~\ref{def:scale-seq}, $(a_K^{(j)})_{K\in 2^{\N}}$ is a scale-$N_j$ $h^{s_j}$-envelope, with
\[
s_j=
\begin{cases}
	s,& v_{N_j}^{(\ast_j)} \text{ is of type (D)},\\
	s_C:=\alpha-\frac1q-\delta,& v_{N_j}^{(\ast_j)} \text{ is of type (C)}.
\end{cases}
\]
Moreover, since $v_{N_5}^{(\ast_5)}$ is of type~(D) in the first inequality of
Proposition~\ref{prop:quintilinear}, the sequence $(b_M)_{M\in 2^{\N}}$ is a scale-$N_5$ $h^{s_5}$-envelope, where
\[
s_5=\begin{cases}
	s,& v_{N_5}^{(\ast_5)} \text{ is of type (D)},\\
	s_D:=\alpha-\frac12-\delta,& v_{N_5}^{(\ast_5)} \text{ is of type (C)}.
\end{cases}
\]
For $\vec{K}=(K_1,K_2,K_3,K_4)$, set
\[
G_{M,\vec{K},n}(t):=\mathbf{P}_M\mathcal{N}_n^{(4)}
\big(
\mathbf{P}_{K_1}v_{N_1}^{(\ast_1)},
\mathbf{P}_{K_2}v_{N_2}^{(\ast_2)},
\mathbf{P}_{K_3}v_{N_3}^{(\ast_3)},
\mathbf{P}_{K_4}v_{N_4}^{(\ast_4)}
\big)(t).
\]
By the uniform estimate~\eqref{cor:4-linear-1-uniform} or Corollary~\ref{cor:4-linearprob} (see Remark~\ref{rmk:numerologyquintic}),
for every $n\sim M$,
\begin{equation}\label{eq:blockwise-pointwise-case2}
	\|\chi_1(t) G_{M,\vec{K},n}(t)\|_{L_t^\infty}
	\lesssim
	\Lambda\,\frac{K\wedge M}{K}\,
	K^{-\nu_0}\,
	\prod_{j=1}^4 K_j^{s_j}a_{K_j}^{(j)},
\end{equation}
with the same $\Lambda,\nu_0$ as in~\eqref{eq:Lambdanu_0}.

Plugging into~\eqref{eq:Q}, we obtain
\begin{align*}
	\mathcal{Q}&\lesssim \Lambda\Big(\sum_{M}b_M^2\Big(\sum_{\vec{K}}\frac{K\wedge M}{K}\, K^{-\nu_0}\cdot\prod_{j=1}^4 K_j^{s_j}a_{K_j}^{(j)}
	\Big)^2\Big)^{1/2}\\
	&\lesssim \Lambda
	\Big(\sum_{M}b_M^2\,\mathrm{I}_M^2\Big)^{1/2}+\Lambda
	\Big(\sum_{M}b_M^2\,\mathrm{II}_M^2\Big)^{1/2},
\end{align*}
where
\[
\mathrm{I}_M:=\sum_{\substack{K_1,K_2,K_3,K_4\\
		K\leq M
}} K^{-\nu_0}\prod_{j=1}^4 K_j^{s_j}a_{K_j}^{(j)},\qquad \mathrm{II}_M:=\sum_{\substack{K_1,K_2,K_3,K_4\\
		K> M
}} M\, K^{-\nu_0-1}\prod_{j=1}^4 K_j^{s_j}a_{K_j}^{(j)}.
\]

Our goal is to show
\begin{equation}\label{eq:target-Q}
	\mathcal{Q}\lesssim N_{(1)}^{-s}N_{(2)}^{-\delta}.
\end{equation}
The key numerology is
\begin{equation}\label{eq:numerology}
	\nu_0 + s_D - s = \nu_0 + \Big(\alpha - \frac{1}{2} - \delta\Big) - s > 10\delta.
\end{equation}

Without loss of generality, assume $N_1 = N_{(1)} = 2N \ge N_5$ and $N_2$ is the second largest among $N_1,N_2,N_3,N_4$, so that $\max(N_2,N_5)=N_{(2)}$. We bound the input sequences using their envelope properties from Definition~\ref{def:scale-seq}:
\begin{align}
	\|a^{(1)}\|_{\ell^2} &\lesssim N_1^{-s_1}, \label{eq:env1} \\
	\|K_2^{s_2-\delta} a^{(2)}\|_{\ell^2} &\lesssim N_2^{-\delta}, \label{eq:env2} \\
	\|K_j^{s_j} a^{(j)}\|_{\ell^2} &\lesssim 1, \quad j=3,4. \label{eq:env34}
\end{align}
For fixed $K=\max(K_1,K_2,K_3,K_4)$, define
\[
H_K := \sum_{\max_j K_j=K} \prod_{j=1}^4 K_j^{s_j} a_{K_j}^{(j)}.
\]
By isolating the index $i \in \{1,2,3,4\}$ that attains $K_i = K$,
\begin{equation}\label{eq:HK-decomp}
	H_K \le \sum_{i=1}^4 K^{s_i} a_K^{(i)} \prod_{j \neq i} \Big(\sum_{K_j \le K} K_j^{s_j} a_{K_j}^{(j)}\Big).
\end{equation}
By Cauchy--Schwarz and~\eqref{eq:env1}--\eqref{eq:env34},
\begin{align}
	\sum_{K_1 \le K} K_1^{s_1} a_{K_1}^{(1)} &\lesssim K^{s_1} N_1^{-s_1}, \label{eq:sumK1} \\
	\sum_{K_2 \le K} K_2^{s_2} a_{K_2}^{(2)} &\lesssim K^{\delta} N_2^{-\delta}, \label{eq:sumK2} \\
	\sum_{K_j \le K} K_j^{s_j} a_{K_j}^{(j)} &\lesssim_{\epsilon} K^{\epsilon},\quad j=3,4, \label{eq:sumK34}
\end{align}
for any $\epsilon > 0$. Substituting~\eqref{eq:sumK1}--\eqref{eq:sumK34} into~\eqref{eq:HK-decomp} for each choice of the maximum index~$i$, we obtain
\begin{align}
	H_K &\lesssim K^{s_1+\delta+2\epsilon} N_2^{-\delta} a_K^{(1)} \notag \\
	&\quad + K^{s_1+\delta+2\epsilon} N_1^{-s_1} \big(K^{s_2-\delta} a_K^{(2)}\big) \notag \\
	&\quad + K^{s_1+\delta+\epsilon} N_1^{-s_1} N_2^{-\delta} \big(K^{s_3} a_K^{(3)} + K^{s_4} a_K^{(4)}\big). \label{eq:HK-bound}
\end{align}
We can write
\begin{equation}\label{eq:HK-final}
	H_K \lesssim K^{s_1+\delta+2\epsilon} \widetilde{A}_K,
\end{equation}
where
\[
\widetilde{A}_K := N_2^{-\delta} a_K^{(1)} + N_1^{-s_1} \big(K^{s_2-\delta} a_K^{(2)}\big) + N_1^{-s_1} N_2^{-\delta} \big(K^{s_3} a_K^{(3)} + K^{s_4} a_K^{(4)}\big),
\]
and
\begin{equation}\label{eq:Atilde-norm}
	\|\widetilde{A}\|_{\ell_K^2} \lesssim N_1^{-s_1} N_2^{-\delta}.
\end{equation}

We now bound $\mathrm{I}_M$ and $\mathrm{II}_M$. For $\mathrm{I}_M$, by Cauchy--Schwarz in~$K$:
\begin{align}
	\mathrm{I}_M &= \sum_{K \le M} K^{-\nu_0} H_K \lesssim \sum_{K \le M} K^{s_1 - \nu_0 + \delta + 2\epsilon} \widetilde{A}_K \notag \\
	&\le \Big(\sum_{K \le M} K^{2(s_1 - \nu_0 + \delta + 2\epsilon)}\Big)^{1/2} \|\widetilde{A}\|_{\ell_{K}^2} \notag \\
	&\lesssim M^{s_1 - \nu_0 + \delta + 2\epsilon} N_1^{-s_1} N_2^{-\delta}, \label{eq:IM-bound}
\end{align}
where the geometric sum is dominated by $K=M$ since $s_1 - \nu_0 \geq 0$.

Similarly, for $\mathrm{II}_M$:
\[
\mathrm{II}_M = M \sum_{K > M} K^{-\nu_0-1} H_K \lesssim M \sum_{K > M} K^{s_1 - \nu_0 - 1 + \delta + 2\epsilon} \widetilde{A}_K.
\]
Since $s_1 \le 1$ and $\nu_0 > 10\delta$, the exponent $s_1 - \nu_0 - 1 + \delta + 2\epsilon$ is strictly negative for small $\epsilon$. By Cauchy--Schwarz,
\begin{align}
	\mathrm{II}_M &\le M \Big(\sum_{K > M} K^{2(s_1 - \nu_0 - 1 + \delta + 2\epsilon)}\Big)^{1/2} \|\widetilde{A}\|_{\ell_{K}^2} \notag \\
	&\lesssim M^{s_1 - \nu_0 + \delta + 2\epsilon} N_1^{-s_1} N_2^{-\delta}. \label{eq:IIM-bound}
\end{align}
The bounds for $\mathrm{I}_M$ and $\mathrm{II}_M$ coincide. Setting $\rho := s_1 - \nu_0 + \delta + 2\epsilon$, we obtain
\begin{equation}\label{eq:Q-intermediate}
	\mathcal{Q} \lesssim \Lambda\, N_1^{-s_1} N_2^{-\delta} \Big(\sum_M b_M^2 M^{2\rho}\Big)^{1/2}.
\end{equation}

We perform the summation over the output scale~$M$ by distinguishing two subcases.

\smallskip\noindent
\textbf{Subcase 2.1: $v_{N_5}^{(\ast_5)}$ is of type~(C).}
The scale is fixed at $M = N_5$, so there is no dyadic sum over~$M$. We have $b_M = \|\mathbf{P}_{N_5} v_{N_5}^{(\mathrm{C})}\|_{X^{0,b}} \lesssim N_5^{-s_D}$. From~\eqref{eq:Q-intermediate},
\[
\mathcal{Q} \lesssim \Lambda\, N_5^{-s_D+\rho} N_1^{-s_1} N_2^{-\delta} = \Lambda\, N_5^{-s_D + s_1 - \nu_0 + \delta + 2\epsilon} N_1^{-s} N_1^{-(s_1 - s)} N_2^{-\delta}.
\]
Since $M = N_5 \le N_1$ and $s_1 \ge s$, we have $N_1^{-(s_1 - s)} \le N_5^{-(s_1 - s)}$. The total power of~$N_5$ is then $-s_D - \nu_0 + s + \delta + 2\epsilon$.
By~\eqref{eq:numerology}, $\nu_0 + s_D - s > 10\delta$. Choosing $\epsilon$ small enough ensures this exponent is strictly negative, yielding $\mathcal{Q} \lesssim \Lambda\, N_{(1)}^{-s} N_{(2)}^{-\delta}$.

\smallskip\noindent
\textbf{Subcase 2.2: $v_{N_5}^{(\ast_5)}$ is of type~(D).}
In this case, $(b_M)_{M \le N_5}$ is a scale-$N_5$ $h^s$-envelope, and we must sum over all dyadic~$M$. We bound
\begin{equation}\label{eq:sum-bM}
	\Big(\sum_{M} M^{2\rho} b_M^2\Big)^{1/2}.
\end{equation}
We further distinguish the type of $v_{N_1}^{(\ast_1)}$. If $v_{N_1}^{(\ast_1)}$ is of type~(D), then $s_1=s$ and $\rho=s-\nu_0+\delta+2\epsilon<s$. By property~(2) of Definition~\ref{def:scale-seq}, the expression~\eqref{eq:sum-bM} is bounded by $N_5^{-(s-\rho)}$.

If $v_{N_1}^{(\ast_1)}$ is of type~(C), then $s_1=s_C$ and $\rho=s_C-\nu_0+\delta+2\epsilon$. If $\rho\leq s$, we apply property~(2) of Definition~\ref{def:scale-seq} with $s_0=\rho$ to bound~\eqref{eq:sum-bM} by $N_5^{-(s-\rho)}\leq 1$. Inserting this into~\eqref{eq:Q-intermediate} yields
\[
\mathcal{Q}\lesssim \Lambda\, N_1^{-s_C}N_2^{-\delta}\leq \Lambda\, N_1^{-s}N_{(2)}^{-\delta},
\]
since $s_C>s+\delta$.

If $\rho>s$, we bound~\eqref{eq:sum-bM} by
\[
\Big(\sum_{M\leq N_5}M^{2\rho}b_M^2\Big)^{1/2}+\Big(\sum_{M>N_5}M^{2\rho}b_M^2\Big)^{1/2}.
\]
The first term is bounded by $N_5^{\rho-s}$. For the second, property~(3) of Definition~\ref{def:scale-seq} gives $|b_M|\leq (M/N_5)^{-1}\cdot N_5^{-s}$, hence
\[
\Big(\sum_{M>N_5}M^{2\rho} b_M^2\Big)^{1/2}\lesssim N_5^{-s}\cdot N_5^{\rho-s},
\]
where we used $\rho<s_C<1$. Inserting into~\eqref{eq:Q-intermediate} gives
\[
\mathcal{Q}\lesssim \Lambda\, N_1^{-s_C}N_2^{-\delta}N_5^{\rho-s}\leq \Lambda\, N_1^{-s}N_2^{-\delta}N_5^{\rho-s}\cdot N_1^{-(s_C-s)}\leq \Lambda\, N_1^{-s}N_2^{-\delta}N_5^{\rho-s_C}.
\]
Since $\rho-s_C<-\delta$, we finally obtain
\[
\mathcal{Q}\lesssim \Lambda\, N_{(1)}^{-s}N_{(2)}^{-\delta}.
\]
This completes the proof of Case~2 and of Proposition~\ref{prop:quintilinear}.
\end{proof}

\end{document}
